\lecture{15}
\title{Short-Time Fourier Transform}

\begin{document}

\subsection{Lecture: Spectrograms}

If you were to shove a recording of a piano playing a single chord
through the DFT, you'd be able to retrieve the frequencies of that chord.
But if you were to shove a recording of a piano playing an entire song
through the DFT, you'd only retrieve the set of all frequencies
present in the entire song---not the timing of each frequency,
or even what order they came in.

The DFT no longer works when \emph{frequencies change over time}.
If we want to analyze the frequencies in an entire song,
we need data that shows the frequencies present at each moment in time.

\begin{center}
    \frame{\includegraphics[width=0.4\textwidth]{lecture-15/spectrogram}}
\end{center}

Such data is called a \textbf{spectrogram}, pictured above.
It is produced by performing several short-time DFTs.

\begin{center}
    \frame{\includegraphics[width=0.4\textwidth]{lecture-15/spectrogram-construction}}
\end{center}

\begin{problem}
    You're given a DT signal with sampling rate
    $f_s = 44100$ samples per second and want a spectrogram such that:
    \begin{itemize}
        \item Sixteenth notes at a tempo of $120$ beats per minute
        (each lasting $\frac{1}{8}$ of a second)
        are distinguished from each other.
        \item Errors in frequency are less than or equal to $5$ Hz.
    \end{itemize}
    Determine the range of all possible window lengths $N$ you might use.
\end{problem}

\begin{solution}
    The first condition implies $\frac{N}{f_s} \leq \frac{1}{8}$,
    and the second condition implies $\frac{f_s}{N} \leq 10$.
    Together, this implies $N \in [4410, 5512]$.
\end{solution}

This process of taking DFTs over sliding windows
is known as the \textbf{Short-Time Fourier Transform (STFT)}.

\begin{remark}
    Later on, $N$ will change from denoting the window size
    to denoting the DFT size.  Those two will become different!
\end{remark}

\subsection{Lecture: Streaming Windows}

Recall from \href{10.html}{Lecture 10} that
filtering a DT signal $x[n]$
is equivalent to multiplying its DTFT $X(\Omega)$ by
a window-shaped frequency response $H(\Omega)$,
which is itself equivalent to convolving
$x[n]$ with the unit-sample response $h[n]$.

In the following discussion,
we'll try to isolate the melody of the following.

\includeaudio[width=0.4\textwidth]{lecture-15/original}

When working with finite-length DFTs,
two issues arise.
\begin{enumerate}
    \item The FFT runs in $O(N \log N)$ time.
    For very long DT signals $x[n]$,
    it is not feasible to filter the entire signal at once.
    \item Filtering DT signals via \emph{circular} convolution
    gives rise to \emph{time aliasing}
    (see \href{13.html#lecture-and-recitation-circular-convolution}{Lecture 13}),
    yielding buggy noises.
\end{enumerate}

A natural fix for the first issue
is to apply STFT:

\begin{quote}
    \textbf{Bad Approach.}
    Break the signal into windows of length $N$,
    filter each window separately,
    and concatenate the results.
    \frame{\includegraphics[width=0.4\textwidth]{lecture-15/circular-filter}}
    Here's how the result sounds:
    \includeaudio[width=0.4\textwidth]{lecture-15/circular-filter}
\end{quote}

Clearly something sounds off.
The main issue comes from our neglect of issue \#2:
circular convolution and time-aliasing.

\begin{remark}
    It might seem that circular convolution
    has no flaws.  Here's an incorrect argument as to why.
    \begin{quote}
        \textbf{False Argument.}
        Let $h[n]$ denote the DT signal used as a filter---the
        inverse DFT of a filtering frequency response.
        From \href{13.html#lecture-and-recitation-circular-convolution}{Lecture 13},
        the circular convolution $y[n] := (x \circledast h)[n]$
        is a DT signal whose DFT only contains good frequencies.

        Therefore, the DTFT concatenation of all of the $y[n]$
        must only contain good frequencies, too.
    \end{quote}
    Here's the flaw in the above: consider the following two extensions on $y[n]$
    over the domain $\mathbb{N}$, instead of just $[0, N - 1]$.
    \[
    y_1[n] := y[n \% N]
    ~~~ \text{ and } ~~~
    y_2[n] := \begin{cases}
    y[n] & \text{ if } n \in [0, N - 1]. \\
    0 & \text{ otherwise.}
    \end{cases}
    \]
    The DFT of $y[n]$ does only contain good frequencies,
    but the implication of this is that the DTFT of $y_1[n]$
    only contains good frequencies---not $y_2[n]$.
    Yet the concatenation of all of the $y[n]$
    is the sum of the $y_2[n]$---not the $y_1[n]$.

    Therefore, there is no reason for
    the DTFT of the concatenation of the $y[n]$
    to only contain good frequencies.
\end{remark}

We really should be taking a linear convolution,
not a circular convolution,
and allowing for time in one window to overflow into the next,
rather than aliasing it backwards.

\begin{quote}
    \textbf{Good Approach (Overlap-Add).}
    Break the signal into windows of length $s$.
    Filter each window separately
    by computing its \emph{linear} convolution
    with a filter $h[n]$ of length $N - s + 1$ (or less).
    Add together the results.

    \frame{\includegraphics[width=0.4\textwidth]{lecture-15/overlap-add}}

    Note that the result of each window's convolution
    may overlap with the results from other windows---hence
    the name ``overlap-add'' and not ``overlap-concatenate''.
\end{quote}

From this point on, it's important to think
less in the time domain $h[n]$,
and more in the frequency domain $H[k]$.
We can rephrase overlap-add
from the perspective of the frequency domain
like so:

\begin{quote}
    \textbf{Good Approach (Overlap-Add --- Frequency Domain).}
    Break $x[n]$ into subsignals $x_i[n]$ of length $s$.
    For each subsignal, compute its length-$N$ DFT $X_i[k]$
    by zero-padding $x_i[n]$ to have length $N$.

    Now compute the element-wise product $X_i'[k] := X_i[k] \cdot H[k]$,
    for our choice of length-$N$ filtering frequency response $H[k]$.
    It is because the filter $h = \mathrm{DFT}^{-1}(H)$ has length $\leq N - s + 1$
    that this filter bypasses any time-aliasing problems.

    Finally, compute the length-$N$ subsignals
    $x'_i := \mathrm{DFT}^{-1}(X_i')$ and sum them together.
\end{quote}

In summary, choosing $H[k]$ is enough
to define the map from length-$s$ subsignals $x_i[n]$
to length-$N$ subsignals $x'_i[n]$.

\begin{remark}
    The frequency domain is how we would go about
    algorithmically implementing this filter, too!
    Naive linear convolution is $O(N^2)$,
    but two DFT's plus an element-wise product is $O(N \log N)$.
\end{remark}

Thus, it remains to construct a length-$N$ frequency response $H[k]$
whose unit-sample response $h[n]$ has length $\leq N - s + 1$.


\subsection{Lecture: Short-Time Filter Design}

If we pick $H[k]$ to act as a perfect filter,
then $h[n]$ will have length $N$.
We must compromise in our design of $H[k]$
so that the last $s$ entries of $h[n]$ are all zeroes.
Here's how.

\begin{quote}
    \textbf{Box Approach.}
    Let's work in reverse: we'll pick $h[n]$
    so its last $s$ entries are all zeroes, then hope $H[k]$
    is an effective filter.
    \begin{enumerate}
        \item Take a perfect (infinite, continuous)
        filtering frequency response $H_1(\Omega)$.  No compromises here.
        \item Define $h_1[n]$ to be the infinite impulse response
        $\mathrm{DTFT}^{-1}(H_1(\Omega))$, \emph{truncated} to length $N_f = N - s + 1$.
        \begin{center}
            \frame{\includegraphics[width=0.4\textwidth]{lecture-15/DT-h1}}
        \end{center}
        \item Define $h_2[n]$ to be a version of $h_1[n]$
        that is zero-padded to have length $N$.
        \begin{center}
            \frame{\includegraphics[width=0.4\textwidth]{lecture-15/DT-h2}}
        \end{center}
        \item Define $H_2[k] := \mathrm{DFT}(h_2[n])$
        of length $N$.
        \begin{center}
            \frame{\includegraphics[width=0.4\textwidth]{lecture-15/H2}}
        \end{center}
    \end{enumerate}
    Then $H_2[k]$ is our designed frequency response $H[k]$.
    Here's how the result sounds:
    \includeaudio[width=0.4\textwidth]{lecture-15/box}
\end{quote}

That works pretty well!  But if you listen closely,
you can still faintly hear some of the frequencies from the harmony,
especially in the second half.
You can see why in the graph of $H_2[k]$:
the ``ripples'' mean it's not a perfect filter.

From what compromise do these ripples come?
The answer is \emph{Step \#2};
ideally, we would take $h_1[n] := \mathrm{DTFT}^{-1}(H_1(\Omega))$
to be an \emph{infinite} DT signal, but instead,
we must compromise by \emph{truncating} to have length $N - s + 1$.

Here's the trick, though:
what if we \emph{truncated} a bit more smoothly?

\begin{quote}
    \textbf{Hann Window Approach.}
    We take the same approach as before,
    but we instead define $h_1[n]$ as follows.
    \[ h_1[n] := w_h[n] \cdot \mathrm{DTFT}^{-1}(H_1(\Omega))
    ~~~ \text{ where } ~~~
    w_h[n] := \begin{cases}
    \sin\left(\frac{\pi n}{N_f}\right)^2 & \text{ if } n \in [0, N_f - 1] \\
    0 & \text{ otherwise.}
    \end{cases} \]
    The signal $w_h[n]$ is called a \textbf{Hann Window},
    in contrast to the typical rectangular window
    \[ w_r[n] := \begin{cases}
    1 & \text{ if } n \in [0, N_f - 1] \\
    0 & \text{ otherwise.}
    \end{cases} \]
    used in the previous approach.
    An alternative choice is a triangle window;
    you can observe the efficacy of each of these truncation methods
    through the diagram below.
    \begin{center}
        \frame{\includegraphics[width=0.4\textwidth]{lecture-15/windows}}
    \end{center}
    In the previous approach,
    the frequencies present in $h_1[n]$
    would be the result of convolving our perfect filter $H_1(\Omega)$
    with the ripple-filled signal $\mathrm{DTFT}(w_r[n])$.
    No wonder the result $H_2[k]$ had so many ripples!

    Expectedly, using the Hann window to construct the truncation $h_1[n]$
    gives a much cleaner $H_2[k]$.
    \begin{center}
        \frame{\includegraphics[width=0.4\textwidth]{lecture-15/Hann-window}}
    \end{center}
    And it sounds much better, too---no more
    harmonies still lingering around.
    \includeaudio[width=0.4\textwidth]{lecture-15/hann}
\end{quote}

\subsection{Recitation: Windows and Decibels}

We already saw how the rectangular, triangular, and Hann windows compare.
But let's compare them at a logarithmic scale.

\begin{center}
    \frame{\includegraphics[width=0.4\textwidth]{lecture-15/window-decibel}}
\end{center}

In the above, we define decibels at a \emph{relative} scale.
(Of course we do---turning the volume knob up shouldn't alter these graphs.)

\textbf{Definition.}
Given a signal $x[n]$, say its highest value is $M$.
Then $x[n]$ in decibels is $10 \log_{10}\left(\frac{x[n]}{M}\right)$.

So a measurement of $-20$ dB corresponds to
a measurement in the signal that is $\times \frac{1}{100}$ the highest measurement.
For reference, measurements of $-60$ dB are essentially inaudible,
whereas measurements of $\approx -25$ dB are still clearly perceptible.

\begin{remark}
    The relationship between windows and their Fourier transforms
    is actually what Heisenberg's uncertainty principle is about!
    Roughly speaking, it says that a tightly-localized frequency response
    cannot correspond to a tightly-localized signal.
\end{remark}

\subsection{Recitation: Short-Time Fourier Transform}

We never actually formally defined the STFT.  Let's do that.

\textbf{Definition.} Given a signal $x[n]$ and a window $w[n]$,
the \textbf{STFT} is the two-dimensional array:
\[ X[k, m] := \frac{1}{N} \sum_{n = 0}^{N - 1}
x[n + ms] \cdot w[n] \cdot e^{-jk \frac{2\pi}{N} n}. \]
Note that $k$ and $m$ each only appear once on the RHS.
The value of $k$ controls the frequency being measured,
and the value of $m$ controls the time being measured.
A \textbf{spectrogram} is a diagram that displays $|X[k, m]|^2$.

Note that the STFT of $x[n]$ depends on three different parameters:
\[ w[n]: \text{ the type of window } ~~~ \parallel ~~~
N: \text{ the DFT size } ~~~ \parallel ~~~
s: \text{ the step size } \]
Importantly, these parameters do not have to be equal!
\begin{itemize}
    \item You can take $s < \mathrm{len}(w[n])$,
    meaning each data point will be used multiple times!
    \begin{itemize}
        \item Or $s > \mathrm{len}(w[n])$, meaning
        some data points are completely ignored\dots
    \end{itemize}
    \item You can take $N > \mathrm{len}(w[n])$,
    meaning the data collected from the window $w[n]$
    will be zero-padded to have length $N$.
    \begin{itemize}
        \item It does not make sense to take $N < \mathrm{len}(w[n])$.
        You can not take a DFT of size $6$ on a DT signal of length $7$.
    \end{itemize}
\end{itemize}
Compare the cleanliness of a rectangular window (at left)
versus a Hann window (at right).
\begin{center}
    \frame{\includegraphics[width=0.6\textwidth]{lecture-15/spectrogram-windows}}
\end{center}
Expectedly, the parameter settings affect the spectrogram as follows.
\begin{itemize}
    \item Increasing $\mathrm{len}(w[n])$ increases the frequency resolution.
    \item Increasing $N$ increases the fineness of frequency measurements---but
    only superficially.  (Recall \href{13.html#lecture-frequency-resolution}{Lecture 13}.)
    \item Decreasing $s$ increases the fineness of time measurements.
\end{itemize}
Of course, these improvements to parameter settings
come at a cost of time complexity.

\end{document}